\section{Экспериментальная оценка реализации}
\label{sec:Chapter7} \index{Chapter7}

\subsection{Цели и методика оценки}
\label{sec:eval_goals}

В рамках оценки решаются следующие задачи:

\begin{itemize}
    \item измерение ускорения времени выполнения программ, получаемого за счёт применения предложенной оптимизации;
    \item оценка накладных расходов на этапе компиляции;
\end{itemize}

Экспериментальная методика основана на сравнении двух конфигураций компиляции:
\begin{enumerate}
    \item стандартный компилятор Go;
    \item модифицированная версия компилятора с реализованным проходом SLP-векторизации.
\end{enumerate}

Для измерения времени выполнения использовался встроенный пакет тестирования~\cite{go_testing}, 
а для комплексных нагрузок применялись внешние тестовые наборы. 
Каждый эксперимент выполнялся многократно, после чего применялась утилита \texttt{benchstat}~\cite{benchstat}, 
предназначенная для статистической обработки результатов. 
Данная утилита позволяет агрегировать результаты многократных запусков, вычислять медианные значения, 
рассчитывать погрешность и оценивать статистическую значимость различий.

\subsubsection{Окружение}
Измерения проводились на изолированном сервере, специально предназначенном для замеров производительности. 
Это позволяет уменьшить влияние фоновой активности на результаты измерений.

Использовались следующие тестовые платформы:
\begin{itemize}
    \item x86-64: \texttt{Intel Xeon Gold 6230N};
    \item ARM64: \texttt{HiSilicon Kunpeng 920}.
\end{itemize}

Для тестирования выделялось 4 ядра с помощью команды \texttt{taskset} ~\cite{taskset}.


\subsection{Выбор тестовых наборов}
\label{sec:benchmarks}

Для оценки эффективности реализации был выбран набор тестов, включающий как синтетические микробенчмарки, 
так и реальные приложения.

\subsubsection{Набор синтетических тестов}

Синтетические тесты предназначены для оценки эффективности векторизации в контролируемых условиях. Они включают:
\begin{itemize}
\item арифметические операции над структурами с полями базовых арифмитических типов 
(\texttt{int64}, \texttt{float64}, \texttt{int32}, \texttt{int16}, \texttt{int8});
\item копирование массивов и структур;
\end{itemize}

Каждый тест представляет собой последовательность однотипных независимых операций, 
потенциально пригодных для SIMD-векторизации. 
Такой набор позволяет изолированно оценить качество поиска векторных шаблонов, 
устойчивость к различным типам данных и влияние особенностей 
SSA-представления на генерацию SIMD-инструкций.

Исходный код бенчмарков приведен в месте с кодом векторизатора \cite{go_cl_783220}.

\subsubsection{Макробенчмарки}

Для комплексной оценки применялся набор \texttt{Sweet benchmarks}~\cite{sweet_bench}, 
используемый разработчиками Go для оценки влияния изменений компилятора и рантайма на производительность реальных приложений.
Этот набор включает как набор известных приложений, так и сценарии измерения времени сборки.
Это позволяет оценить эффект на реальных приложениях.

\subsection{Результаты измерения производительности}
\label{sec:perf_results}

\subsubsection{Синтетические тесты}

На микробенчмарках наблюдается наибольший эффект, 
поскольку структура вычислений специально подобрана под предположения алгоритма SLP-векторизации.

\begin{table}[H]
    \centering
    \caption{Результаты синтетических тестов на HiSilicon Kunpeng-920}
    \small
    \begin{tabular}{lcccc}
        \hline
        Бенчмарк & \texttt{base} (sec/op) & \texttt{slp} (sec/op) & Изменение &\\
        \hline
        Vec/Copy/int64-2                 & 5.109µ ± 21\% & 3.656µ ± 1\% & -28.43\%      & (p=0.002) \\
        Vec/Copy/int32-2                 & 5.114µ ± 21\% & 3.656µ ± 0\% & -28.50\%      & (p=0.002) \\
        Vec/Copy/int16-2                 & 3.654µ ± 0\%  & 3.655µ ± 1\% & \texttildelow & (p=0.067) \\
        Vec/Copy/int8-2                  & 3.658µ ± 0\%  & 3.654µ ± 0\% & -0.11\%       & (p=0.002) \\
        Vec/Copy/Unaligned/offset=1-2    & 6.254µ ± 0\%  & 3.798µ ± 7\% & -39.27\%      & (p=0.002) \\
        Vec/Copy/Unaligned/offset=2-2    & 6.253µ ± 0\%  & 3.797µ ± 0\% & -39.28\%      & (p=0.002) \\
        Vec/Copy/Unaligned/offset=3-2    & 6.582µ ± 5\%  & 3.800µ ± 7\% & -42.26\%      & (p=0.002) \\
        Vec/Copy/Unaligned/offset=4-2    & 6.253µ ± 0\%  & 3.803µ ± 0\% & -39.18\%      & (p=0.002) \\
        Vec/Copy/Unaligned/offset=5-2    & 6.253µ ± 0\%  & 3.796µ ± 0\% & -39.29\%      & (p=0.002) \\
        Vec/Copy/Unaligned/offset=6-2    & 6.253µ ± 0\%  & 3.796µ ± 2\% & -39.28\%      & (p=0.002) \\
        Vec/Copy/Unaligned/offset=7-2    & 6.253µ ± 0\%  & 3.798µ ± 0\% & -39.26\%      & (p=0.002) \\
        Vec/Copy/Unaligned/offset=8-2    & 6.251µ ± 0\%  & 3.805µ ± 0\% & -39.13\%      & (p=0.002) \\
        Vec/Copy/Overlap-2               & 13.74µ ± 188\% & 12.25µ ± 9\% & -10.82\%     & (p=0.002) \\
        \hline
        Vec/Simple/Sub/int64-2           & 6.815µ ± 0\%  & 4.753µ ± 1\% & -30.25\%      & (p=0.002) \\
        Vec/Simple/Sub/float64-2         & 6.773µ ± 0\%  & 4.756µ ± 0\% & -29.78\%      & (p=0.002) \\
        Vec/Simple/Sub/int32-2           & 9.176µ ± 0\%  & 4.753µ ± 0\% & -48.20\%      & (p=0.002) \\
        Vec/Simple/Sub/float32-2         & 13.607µ ± 0\% & 4.765µ ± 6\% & -64.98\%      & (p=0.002) \\
        Vec/Simple/Sub/int16-2           & 26.09µ ± 0\%  & 26.15µ ± 0\% & \texttildelow & (p=0.180) \\
        Vec/Simple/Sub/int8-2            & 48.55µ ± 2\%  & 49.23µ ± 1\% & +1.39\%       & (p=0.019) \\
        \hline
        Vec/Complex/int64-2              & 8.561µ ± 0\%  & 6.496µ ± 0\% & -24.13\%      & (p=0.002) \\
        Vec/Complex/float64-2            & 9.603µ ± 0\%  & 6.876µ ± 1\% & -28.40\%      & (p=0.002) \\
        Vec/Complex/int32-2              & 13.479µ ± 0\% & 6.516µ ± 0\% & -51.66\%      & (p=0.002) \\
        Vec/Complex/float32-2            & 17.946µ ± 0\% & 6.920µ ± 2\% & -61.44\%      & (p=0.002) \\
        Vec/Complex/int16-2              & 36.89µ ± 0\%  & 37.42µ ± 2\% & \texttildelow & (p=0.065) \\
        Vec/Complex/int8-2               & 64.10µ ± 0\%  & 64.05µ ± 0\% & \texttildelow & (p=0.394) \\
        \hline
        \textit{geomean}                 & 9.645µ & 6.595µ & -31.62\% \\
        \hline
    \end{tabular}
\end{table}

\begin{table}[H]
    \centering
    \caption{Результаты синтетических тестов на Intel Xeon Gold 6230N CPU @ 2.30GHz}
    \small
    \begin{tabular}{lcccc}
        \hline
        Бенчмарк & \texttt{base} (sec/op) & \texttt{slp} (sec/op) & Изменение &\\
        \hline
        Vec/Copy/int64-2                 & 3.602µ ± 0\% & 3.390µ ± 5\% & -5.89\%       & (p=0.002) \\
        Vec/Copy/int32-2                 & 3.608µ ± 0\% & 3.330µ ± 1\% & -7.71\%       & (p=0.002) \\
        Vec/Copy/int16-2                 & 3.319µ ± 0\% & 3.326µ ± 1\% & \texttildelow & (p=0.013) \\
        Vec/Copy/int8-2                  & 3.320µ ± 0\% & 3.322µ ± 1\% & \texttildelow & (p=0.015) \\
        Vec/Copy/Unaligned/offset=1-2    & 4.191µ ± 0\% & 3.485µ ± 1\% & -16.85\%      & (p=0.002) \\
        Vec/Copy/Unaligned/offset=2-2    & 4.185µ ± 0\% & 3.483µ ± 3\% & -16.77\%      & (p=0.002) \\
        Vec/Copy/Unaligned/offset=3-2    & 4.191µ ± 0\% & 3.482µ ± 0\% & -16.91\%      & (p=0.002) \\
        Vec/Copy/Unaligned/offset=4-2    & 4.185µ ± 0\% & 3.482µ ± 0\% & -16.80\%      & (p=0.002) \\
        Vec/Copy/Unaligned/offset=5-2    & 4.185µ ± 0\% & 3.482µ ± 0\% & -16.81\%      & (p=0.002) \\
        Vec/Copy/Unaligned/offset=6-2    & 4.185µ ± 0\% & 3.486µ ± 0\% & -16.71\%      & (p=0.002) \\
        Vec/Copy/Unaligned/offset=7-2    & 4.185µ ± 0\% & 3.485µ ± 0\% & -16.73\%      & (p=0.002) \\
        Vec/Copy/Unaligned/offset=8-2    & 3.699µ ± 3\% & 3.494µ ± 3\% & -5.56\%       & (p=0.002) \\
        Vec/Copy/Overlap-2               & 9.047µ ± 0\% & 6.409µ ± 0\% & -29.15\%      & (p=0.002) \\
        \hline
        Vec/Simple/Sub/int64-2           & 5.445µ ± 0\% & 4.277µ ± 0\% & -21.44\%      & (p=0.002) \\
        Vec/Simple/Sub/float64-2         & 5.837µ ± 1\% & 4.273µ ± 12\% & -26.79\%     & (p=0.002) \\
        Vec/Simple/Sub/int32-2           & 8.810µ ± 1\% & 4.270µ ± 0\% & -51.54\%      & (p=0.002) \\
        Vec/Simple/Sub/float32-2         & 9.575µ ± 0\% & 4.269µ ± 2\% & -55.42\%      & (p=0.002) \\
        Vec/Simple/Sub/int16-2           & 21.66µ ± 0\% & 21.62µ ± 0\% & \texttildelow & (p=0.065) \\
        Vec/Simple/Sub/int8-2            & 35.71µ ± 0\% & 35.75µ ± 0\% & \texttildelow & (p=0.394) \\
        \hline
        Vec/Complex/int64-2              & 6.646µ ± 1\% & 5.024µ ± 0\% & -24.41\%      & (p=0.002) \\
        Vec/Complex/float64-2            & 6.879µ ± 1\% & 5.030µ ± 0\% & -26.88\%      & (p=0.002) \\
        Vec/Complex/int32-2              & 13.231µ ± 0\% & 5.025µ ± 1\% & -62.02\%     & (p=0.002) \\
        Vec/Complex/float32-2            & 11.611µ ± 2\% & 5.035µ ± 2\% & -56.64\%     & (p=0.002) \\
        Vec/Complex/int16-2              & 29.04µ ± 2\% & 29.12µ ± 0\% & \texttildelow & (p=0.394) \\
        Vec/Complex/int8-2               & 48.80µ ± 3\% & 48.97µ ± 0\% & \texttildelow & (p=0.058) \\
        \hline
        \textit{geomean}                 & 7.144µ & 5.552µ & -22.29\% \\
        \hline
    \end{tabular}
\end{table}

Полученные результаты показывают устойчивое ускорение для большинства операций. 

Ключевыми наблюдениями являются:
\begin{itemize}
\item копирование и арифметика над \texttt{int32/float32} дают максимальный эффект (до \textbf{~60\%});
\item 64-битные операции стабильно ускоряются на \textbf{~20--28\%};
\item невыровненные доступы не являются критической проблемой и показывают прирост производительности;
\item не векторизируются типы \texttt{int8} и \texttt{int16};
\end{itemize}

Дополнительно выявлены системные ограничения, влияющие на эффективность SLP-векторизации:

\begin{enumerate}
\item \textbf{Разделение базовых блоков из-за проверок границ обращения к срезам.} \
Проверки выхода за границы срезов разделяют базовые блоки. 
Поскольку SLP-векторизация работает только внутри одного базового блока, 
это препятствует объединению операций.

\item \textbf{Ограничения модели памяти.} \\
SLP может объединять только те операции загрузки, которые имеют \emph{общее состояние памяти}. 
Для этого в тестах приходилось явно поднимать загрузки данных выше.
\begin{itemize}
    \item \textbf{Мелкие типы данных.} \\
    Для \texttt{int8} и \texttt{int16} компилятор генерирует последовательности \texttt{Load--ALU Op--Store},
    где загрузка данных каждой следующей цепочки принимает состояние памяти от предыдущей, 
    что нарушает необходимое условие общего состояния памяти для их объединения.
\end{itemize}

\end{enumerate}


Синтетические тесты подтверждают, что предложенный подход к векторизации обеспечивает 
устойчивое ускорение на широком классе задач, 
особенно для 32- и 64-битных операций. 
Основные ограничения связаны не с качеством поиска SIMD-шаблонов, 
а с архитектурными ограничениями и особенностями SSA-формы: границами базовых блоков, моделью памяти и
невозможностью объединения зависимых операций загрузки.

\subsubsection{Макробенчмарки}

Результаты макробенчмарков, ожидаемо, показывают, 
что предложенный прототип SLP-векторизации пока не оказывает заметного влияния на общую производительность ни на x86, ни на ARM64.

На обеих архитектурах геометрическое среднее остаётся практически неизменным, 
что ожидаемо для первого консервативного прототипа. 

Однако отдельно выделяется бенчмарк \texttt{Tile38QueryLoad} на x86, 
показавший устойчивое улучшение. 

\begin{table}[htbp]
    \centering
    \caption{Результаты \texttt{Sweet} на \texttt{HiSilicon Kunpeng-920}}
    \small
    \begin{tabular}{lcccc}
        \hline
        Бенчмарк & \texttt{base} (sec/op) & \texttt{slp} (sec/op) & Изменение &\\
        \hline
        BleveIndexBatch100-4      & 7.467 ± 1\% & 7.513 ± 2\% & \texttildelow   & (p=0.165) \\
        ESBuildThreeJS-4          & 755.1m ± 1\% & 754.5m ± 1\% & \texttildelow & (p=0.912) \\
        ESBuildRomeTS-4           & 195.5m ± 2\% & 194.0m ± 1\% & \texttildelow & (p=0.123) \\
        EtcdPut-4                 & 55.08m ± 1\% & 54.69m ± 1\% & \texttildelow & (p=0.165) \\
        EtcdSTM-4                 & 292.2m ± 1\% & 291.4m ± 1\% & \texttildelow & (p=0.436) \\
        GopherLuaKNucleotide-4    & 32.99 ± 3\%  & 31.92 ± 2\% & -3.23\%         & (p=0.004) \\
        MarkdownRenderXHTML-4     & 265.8m ± 0\% & 265.0m ± 0\% & -0.28\%       & (p=0.001) \\
        Tile38QueryLoad-4         & 607.9µ ± 0\% & 608.0µ ± 0\% & \texttildelow & (p=0.739) \\
        \hline
        \textit{geomean}          & 310.0m & 308.4m & -0.49\% \\
        \hline
    \end{tabular}
\end{table}

\begin{table}[htbp]
    \centering
    \caption{Результаты \texttt{Sweet} на \texttt{Intel Xeon Gold 6230N}}
    \small
    \begin{tabular}{lcccc}
        \hline
        Бенчмарк & \texttt{base} (sec/op) & \texttt{slp} (sec/op) & Изменение & \\
        \hline
        BleveIndexBatch100-4      & 5.867 ± 2\% & 5.842 ± 2\% & \texttildelow     & (p=0.971) \\
        ESBuildThreeJS-4          & 685.0m ± 2\% & 686.9m ± 1\% & \texttildelow   & (p=0.853) \\
        ESBuildRomeTS-4           & 179.6m ± 1\% & 180.2m ± 1\% & \texttildelow   & (p=0.353) \\
        EtcdPut-4                 & 138.7m ± 37\% & 134.6m ± 19\% & \texttildelow & (p=0.315) \\
        EtcdSTM-4                 & 526.0m ± 4\% & 531.9m ± 2\% & \texttildelow   & (p=0.579) \\
        GopherLuaKNucleotide-4    & 27.62 ± 1\% & 27.63 ± 1\% & \texttildelow     & (p=0.739) \\
        MarkdownRenderXHTML-4     & 237.6m ± 1\% & 239.7m ± 1\% & +0.89\%         & (p=0.001) \\
        Tile38QueryLoad-4         & 532.3µ ± 0\% & 523.3µ ± 0\% & -1.69\%         & (p=0.000) \\
        \hline
        \textit{geomean}          & 341.3m & 340.2m & -0.31\% \\
        \hline
    \end{tabular}
\end{table}

На архитектуре ARM в тесте \texttt{GopherLuaKNucleotide} наблюдается ускорение порядка 3.2\%.
На x86 в тесте \texttt{Tile38QueryLoad} зафиксировано ускорение порядка 1.7\%. 
При этом в бенчмарке \texttt{MarkdownRenderXHTML} наблюдается незначительная регрессия порядка 0.9\%, 
что находится в пределах статистической погрешности (около 1\%).

Таким образом, текущая реализация оказывает локальные улучшения 
на отдельных нагрузках при сохранении стабильной производительности в среднем.

\subsubsection{Влияние на время компиляции}
\label{sec:compile_time}

Накладные расходы на компиляцию оценены с помощью набора \texttt{GoBuild} из \texttt{Sweet}.

\begin{table}[htbp]
    \centering
    \caption{Время сборки \texttt{Sweet/GoBuild} на  \texttt{HiSilicon Kunpeng-920}}
    \small
    \begin{tabular}{lcccl}
        \hline
        Тест & \texttt{base} (sec/op) & \texttt{slp} (sec/op) & Изменение & $p$-value \\
        \hline
        GoBuildKubelet-4       & 158.9 $\pm$ 0\% & 161.9 $\pm$ 0\% & +1.85\%  & 0.000 \\
        GoBuildKubeletLink-4   & 12.60 $\pm$ 1\% & 12.61 $\pm$ 0\% & $\sim$   & 0.631 \\
        GoBuildIstioctl-4      & 124.7 $\pm$ 0\% & 126.6 $\pm$ 0\% & +1.47\%  & 0.000 \\
        GoBuildIstioctlLink-4  & 8.550 $\pm$ 1\% & 8.543 $\pm$ 0\% & $\sim$   & 0.247 \\
        GoBuildFrontend-4      & 46.19 $\pm$ 0\% & 52.08 $\pm$ 1\% & +12.75\% & 0.000 \\
        GoBuildFrontendLink-4  & 2.148 $\pm$ 1\% & 2.158 $\pm$ 1\% & $\sim$   & 0.190 \\
        GoBuildTsgo-4          & 76.48 $\pm$ 1\% & 77.17 $\pm$ 2\% & +0.89\%  & 0.003 \\
        GoBuildTsgoLink-4      & 1.163 $\pm$ 1\% & 1.162 $\pm$ 1\% & $\sim$   & 0.971 \\
        \hline
        \textit{geomean}       & 19.25           & 19.65           & +2.10\%  & \\
        \hline
    \end{tabular}
\end{table}

\begin{table}[htbp]
    \centering
    \caption{Время сборки \texttt{Sweet/GoBuild} на \texttt{Intel Xeon Gold 6230N}}
    \small
    \begin{tabular}{lcccc}
        \hline
        Тест & \texttt{base} (sec/op) & \texttt{slp} (sec/op) & Изменение &\\
        \hline
        GoBuildKubelet-4       & 145.8 ± 1\% & 148.3 ± 1\% & +1.68\%       & (p=0.000) \\
        GoBuildKubeletLink-4   & 9.671 ± 0\% & 9.636 ± 0\% & -0.36\%       & (p=0.023) \\
        GoBuildIstioctl-4      & 118.0 ± 1\% & 119.8 ± 0\% & +1.55\%       & (p=0.000) \\
        GoBuildIstioctlLink-4  & 10.14 ± 0\% & 10.12 ± 0\% & \texttildelow & (p=0.280) \\
        GoBuildFrontend-4      & 41.94 ± 1\% & 44.33 ± 0\% & +5.70\%       & (p=0.000) \\
        GoBuildFrontendLink-4  & 1.631 ± 1\% & 1.634 ± 1\% & \texttildelow & (p=0.971) \\
        GoBuildTsgo-4          & 66.77 ± 0\% & 67.54 ± 1\% & +1.15\%       & (p=0.000) \\
        GoBuildTsgoLink-4      & 856.5m ± 1\% & 845.4m ± 1\% & -1.29\%     & (p=0.004) \\
        \hline
        \textit{geomean}       & 16.88 & 17.06 & +1.04\% \\
        \hline
    \end{tabular}
\end{table}

На x86 фиксируются незначительные регрессии в задачах около 1\%-2\%, 
за исключением \texttt{GoBuildFrontend}, 
который содержит базовый блок с несколькими десятками тысяч инструкций, 
из которых векторизуются более двух тысяч. 
Этим объясняется наблюдаемый прирост времени сборки порядка +5\% на x86 и 12\% на Arm. 

В целом геометрическое среднее увеличения времени компиляции составляет порядка 2\%, 
что является приемлемым для прототипа оптимизации.

\subsection{Обсуждение ограничений и невекторизованных случаев}
\label{sec:limitations}

В ходе экспериментов выявлены сценарии, в которых векторизатор не срабатывает или 
даёт незначительный выигрыш. 
Все они прямо вытекают из компромиссов, перечисленных в разделе~\ref{sec:tradeoffs}.

\begin{itemize}
    \item \textbf{Ограничение 128-битными векторами.}
    Поддержка исключительно 128-битных SIMD-регистров накладывает жёсткое требование: г
    руппа операций должна точно заполнять весь регистр. Если ширина группы не кратна 128\,бит 
    (например, три значения \texttt{int64} или два \texttt{int64} и один \texttt{float64}), 
    векторизация не выполняется, и код остаётся скалярным.

    \item \textbf{Отсутствие модели стоимости.}
    Без оценки накладных расходов векторизатор вынужденно отвергает любые варианты, 
    требующие перестановки элементов между линиями. Это отсекает класс шаблонов, 
    где аргументы операции распределены по разным линиям, даже если перестановка была бы выгодна.

    \item \textbf{Консервативная модель памяти.}
    Требование единого состояния \texttt{Memory} (или строго упорядоченной цепочки для сохранений) 
    является главным источником пропущенных возможностей. В частности, не векторизуются:
    \begin{itemize}
        \item независимые загрузки, разделённые записью по гарантированно непересекающемуся адресу — они получают разные состояния \texttt{Memory} и становятся невидимы для SLP-группировки;
    \end{itemize}
\end{itemize}

Все перечисленные ограничения не являются фундаментальными; они могут быть поэтапно сняты при развитии реализации.

\subsection{Направления будущей работы}
\label{sec:future_work}

С учётом принятых компромиссов и полученных экспериментальных данных дальнейшее развитие 
прототипа целесообразно вести по следующим направлениям:

\begin{itemize}
    \item \textbf{Внедрение упрощённой модели стоимости.}
    Эвристическая оценка количества вставок и извлечений позволит более агрессивно векторизовывать код, 
    в том числе случаи, требующие перестановки элементов, но дающие общий положительный эффект.

    \item \textbf{Улучшение анализа памяти.}
    Разработка более точно анализа, способного распознавать независимость загрузок 
    даже при наличии промежуточных сохранений по непересекающимся адресам, 
    напрямую снимет одно из главных ограничений текущей консервативной модели.

    \item \textbf{Поддержка shuffle-инструкций.}
    Реализация операций перестановки элементов откроет возможность векторизации некогерентных линий, 
    где порядок операндов в разных позициях вектора не совпадает.

    \item \textbf{Расширение ширины векторов.}
    Переход от 128-битных к 256-битным векторам (AVX2 на x86-64, SVE на ARM64) увеличит
    объём одновременно обрабатываемых данных и потенциально повысит эффективность оптимизации для более широких групп операций.

    \item \textbf{Векторизация редукций.}
    Добавление распознавания операций свёртки (суммирование, поиск минимума/максимума и т.п.) 
    позволит применять SLP-векторизацию в ещё более широком классе вычислительных ядер.
\end{itemize}